\section{Phase logic as diagnostic regimes}
\label{sec:pdf2_phase_logic}

\subsection{Purpose and stance}
Phase logic is introduced as a \emph{diagnostic} language for classifying observed
enterprise dynamics into regimes.
It is not a control scheme, not a transformation methodology, and not an
optimization framework.

A \emph{phase} is a formal diagnostic category: a set of internal realizations
(continua or states) sharing invariant properties relevant to feasibility,
stability, and admissible transitions under the given formalization.

\subsection{Phases as subsets of \texorpdfstring{$\Omega(K_{EA})$}{Omega(K EA)} and/or model spaces}

Let $K_{EA}$ be fixed.
Phases may be defined at different formal levels, depending on which internal
structures are made explicit.

\begin{definition}[State-space phase]
\label{def:state_space_phase}
A \emph{state-space phase} is a subset $\Phi \subseteq \Omega(K_{EA})$ specified
by constraints expressed using existing primitives, such as thresholds $\Theta$,
distance to the boundary $\partial\Omega$, admissible flows $J$, existence of
cycles $C$, continuity indicators $k$, and feasibility relations.
\end{definition}

\begin{definition}[Model-space phase]
\label{def:model_space_phase}
A \emph{model-space phase} is a subset
\[
\widehat{\Phi} \subseteq \{(M_{current}, M_{declared}, M_{feasible})\}
\]
specified by constraints on feasibility, consistency, and compatibility between
model spaces.
\end{definition}

\paragraph{Observationally induced phase statement.}
Given observations $D$ collected under $(E,\Pi)$, phase classification is
generally \emph{set-valued} unless diagnosability holds:
\[
\Phi_{admissible}(D) :=
\left\{\Phi \;\middle|\; \exists K \in \mathcal{K}_{compat}(D;E,\Pi)
\text{ such that } K \in \Phi \right\}.
\]
A unique phase label is justified only if it is invariant over the entire
compatibility set.

\subsection{Canonical diagnostic regimes (non-prescriptive)}
We introduce a minimal, canonical family of diagnostic regimes expressed without
adding new primitives.
Each regime is defined through existence or non-existence statements concerning
feasibility, continuity, and boundary proximity.

\begin{definition}[Regime \texorpdfstring{$\mathbf{R_{STOP}}$}{R\_STOP} (feasibility collapse)]
\label{def:regime_stop}
$\mathbf{R_{STOP}}$ holds if $M_{feasible} = \varnothing$, as defined in
Definition~\ref{def:feasibility_collapse_stop}.
This regime represents an epistemic non-existence classification.
\end{definition}

\begin{definition}[Regime \texorpdfstring{$\mathbf{R_{COLLAPSE}}$}{R\_COLLAPSE}]
\label{def:regime_collapse}
$\mathbf{R_{COLLAPSE}}$ holds if the compatibility set implies boundary violation,
in the sense that all compatible realizations lie on or beyond the boundary:
\[
\forall (K \text{ or } x) \text{ compatible with } D:\quad
(K \notin \Omega) \ \text{or}\ x \in \partial\Omega.
\]
Equivalently, if a diagnosable bound enforces $d(\cdot,\partial\Omega)\le 0$.
\end{definition}

\begin{definition}[Regime \texorpdfstring{$\mathbf{R_{INERTIA}}$}{R\_INERTIA} (stable but non-expansive)]
\label{def:regime_inertia}
$\mathbf{R_{INERTIA}}$ holds if:
(i) feasibility is non-empty, and
(ii) under the admissible transition family $\Psi$, only local transitions are
compatible with thresholds, such that no increase in effective structural
dimensionality is diagnostically forced.
\end{definition}

\begin{definition}[Regime \texorpdfstring{$\mathbf{R_{ADAPT}}$}{R\_ADAPT} (local adaptation)]
\label{def:regime_adapt}
$\mathbf{R_{ADAPT}}$ holds if:
(i) feasibility is non-empty, and
(ii) at least one compatible realization admits observed changes as internal
reconfiguration within $\Omega(K_{EA})$ without violating thresholds, with
continuity preserved (e.g.\ non-vanishing $k$ over the considered horizon, if
defined in the core).
\end{definition}

\begin{definition}[Regime \texorpdfstring{$\mathbf{R_{TRANSITION}}$}{R\_TRANSITION} (phase transition admissible)]
\label{def:regime_transition}
$\mathbf{R_{TRANSITION}}$ holds if:
(i) feasibility is non-empty, and
(ii) at least one compatible realization requires or exhibits a transition across
a threshold class that changes admissible structure (e.g.\ activation of a
previously absent axis or cycle family, or a diagnosable change in the admissible
set of $\Psi$-moves), while remaining within $\Omega(K_{EA})$.
\end{definition}

\paragraph{Remarks.}
\begin{itemize}
\item Regimes are defined purely through constraints on existing primitives;
no prescriptive interpretation is implied.
\item The intended use is diagnostic: to state which regimes are supported,
excluded, or undecidable given the evidence.
\item Under partial observability, regimes are not necessarily mutually exclusive.
\end{itemize}

\subsection{Phase logic under underdetermination}
Because $\mathcal{K}_{compat}(D;E,\Pi)$ is typically non-singleton, phase logic is
formulated as a three-valued (or set-valued) diagnostic logic:
\emph{ruled out}, \emph{admissible}, or \emph{forced}.

\begin{definition}[Phase status relative to $D$]
\label{def:phase_status}
Let $\mathbf{R}$ be any regime predicate.
Define:
\begin{align*}
\mathbf{R}\ \textbf{ruled out}
&\iff \forall K \in \mathcal{K}_{compat}(D;E,\Pi):\ \neg \mathbf{R}(K),\\
\mathbf{R}\ \textbf{forced}
&\iff \forall K \in \mathcal{K}_{compat}(D;E,\Pi):\ \mathbf{R}(K),\\
\mathbf{R}\ \textbf{admissible}
&\iff \exists K \in \mathcal{K}_{compat}(D;E,\Pi):\ \mathbf{R}(K).
\end{align*}
\end{definition}

\paragraph{Key consequence.}
A unique phase label is justified only when some regime is forced and all
incompatible regimes are ruled out.
Otherwise, the correct diagnostic output is a set of admissible regimes.

\paragraph{Generic non-uniqueness of phase classification.}
For finite observation traces under partial observability, multiple regimes among
$\{\mathbf{R_{INERTIA}}, \mathbf{R_{ADAPT}}, \mathbf{R_{TRANSITION}}\}$ may remain
admissible simultaneously.
Phase classification is therefore generally not forced.

\paragraph{Interpretation.}
Ambiguous phase outcomes are not a defect of the framework; they reflect
irreducible epistemic limits and prevent unwarranted diagnostic commitment.

\subsection{STOP dominance and logical precedence}
Some regimes dominate others by logical precedence.

\paragraph{STOP dominance.}
If $\mathbf{R_{STOP}}$ is forced by $D$, then no non-STOP regime can be forced.
Any classification asserting a forced non-STOP regime under feasibility collapse
is inconsistent with the definition of STOP.

\subsection{What phase logic does \emph{not} do}
Phase logic does not:
(i) recommend actions;
(ii) rank alternatives;
(iii) encode performance objectives or KPI targets;
(iv) assume controllability or optimizability of $K_{EA}$.

It provides a disciplined language for stating which diagnostic regimes are
supported, excluded, or undecidable under explicit non-identifiability.
